\section{Introduction}\label{sec:intro}

\paragraph{The question.} From its sixth part on, this study has asked whether a learner without continuous feature parameters can do at least as well as gradient-trained networks, with the same training data and no more FLOPs. Its tenth part answered with v3, a stacked portfolio \cite{portfoliolearner2026}. This paper studies one learner instead of an ensemble: the portfolio's Boolean member, run alone. How the object was chosen, and the alternatives set aside, are in \cref{app:choice}.

\paragraph{Background.} Parts 1--5 found that rich-regime networks on Boolean targets grow features greedily, cheapest first, and pay in samples: for a $k$-bit parity on $d$ inputs their $\nstar$, the training-set size at 10\% test error, grows like $d^{1.0}$, $d^{1.9}$ and $d^{2.4}$ for $k=2$, 3 and 4. That is close to the leap and correlational-SQ rate $d^{k-1}$ \cite{abbe2023leap} for $k\le3$, and far from the $\approx k\log d$ of an Occam learner \cite{gdcircuit2026}. Part 6 built DFG from that picture, with enumeration and union-bound tests in place of the networks' noisy search \cite{discretelearners2026}. Parts 8--10 stacked it with classical learners in the portfolios v1--v3, whose toy wins came from DFG, and gave it signed-sum blocks (v1.2), then weighted blocks, forced blocks and a beam search (v1.3) \cite{discretelearners2026,portfoliolearner2026}; \cref{app:refs} records what the earlier versions did.

\paragraph{The object: DFG v1.3, run alone.} DFG (discrete feature growth) is forward stepwise least squares over an explicit, finite Boolean dictionary, with a Bonferroni entry rule per class: in learning-theory terms, an Occam learner made greedy by a residual (\cref{sec:relates}). The dictionary holds Fourier characters $\chi_S(x)=\prod_{i\in S}x_i$, indicators of the values of signed sums $\sum_{i\in S}s_ix_i$ of inputs, and a few composites of these. Each round refits the $\pm1$ labels by least squares on the features found so far and admits every candidate whose correlation with the residual exceeds $s\sqrt{2\ln(4KN/\delta)/n}$ for a class of $N$ candidates, $n$ rows and residual scale $s$, with $K=4$ and $\delta=0.1$; the loop stops when a round admits nothing or every training label's sign is right, and DFG predicts with the sign of its readout (\cref{sec:oneline}). This paper studies version 1.3, frozen in Part 10, run by itself and unchanged (\cref{sec:part11}). It is one thing in a specific sense: one residual loop, one admission principle (a tail-bound threshold for each explicit class, computed with four statistics) and one linear readout. Around it sit a multi-start search for signed-sum blocks, a forced guess at $n\le512$ chosen by cross-validation, a logistic refinement and a Hebbian background. \Cref{sec:ablations} shows which of these can be removed, and \cref{sec:core} runs, post hoc, a core without most of them. The price is the domain: DFG v1.3 takes $\pm1$ inputs and two classes, so it runs on the study's 11 Boolean suites, on 2 of its 26 tabular datasets and on no images.

\paragraph{Why it is worth studying.}
\begin{itemize}
\item \emph{It is analysable.} In round 0, when its degree is searched on all $n$ rows, a $k$-parity passes its test exactly when $n(1-\bar y^2)>2\ln(4K\binom dk/\delta)$, so its sample cost grows like $2k\ln d$, where the networks' grows like $d^{k-1}$ for $k\le3$ (\cref{sec:guarantees}).
\item \emph{It realizes the networks' greedy residual loop with the search's sampling noise replaced by enumeration}, one end of the trade between data and compute (\cref{sec:disc-why}), and \emph{its main failure reduces to one question}: when does a search seeded by degree-1 and degree-2 residual statistics find a signed-sum block (\cref{sec:disc-core})?
\item \emph{It is the part of v3 that is new and that matters on Boolean inputs.} v3's other members are LightGBM \cite{ke2017lightgbm}, ExtraTrees \cite{geurts2006extremely}, nearest-neighbour votes and logistic regression, combined by a convex stacker \cite{wolpert1992stacked}, and its image members LightGBM on fixed banks of pixel-subset sums; Part 8's post-hoc ablations traced the portfolio's toy wins to DFG's character search \cite{discretelearners2026,portfoliolearner2026}. And its fitted model is an explicit list of named features (\cref{sec:interp}).
\end{itemize}

\begin{table}[tbp]
\centering\footnotesize\renewcommand{\arraystretch}{1.0}
\begin{tabularx}{\linewidth}{@{}>{\raggedright\arraybackslash}p{2.2cm}>{\raggedright\arraybackslash}p{2.0cm}>{\raggedright\arraybackslash}X>{\raggedright\arraybackslash}p{4.0cm}@{}}
\toprule
evaluation & reference & DFG v1.3 alone & limits\\
\midrule
held out: \code{dlheld7}, 12 targets (\cref{sec:res-dlheld7}) & \code{nets_best}: best of 3 variants (hindsight) & mean $\log_2\nstar$ ratio $-1.29$ over the 11 targets both learn ($2.4\times$ less data; median $-1.16$; s.d.\ 0.91, s.e.\ 0.27); clearly better on 10 of 12, worse on 0 & one family; on 5 of the 10 wins \code{lin} or \code{hgb} also beats \code{nets_best}; v3 needs less data ($+0.21$; clearly on the 3 controls); no learner learns f7\_05\\
\addlinespace
pooled: 126 targets of 10 suites (\cref{sec:res-suites}) & \code{nets_best}: best of 2 to 5 variants & $-1.87$ over 123 ($3.7\times$); clearly better on 117, worse on 2 & in-sample (\cref{tab:status}); against \code{nets_wd} the margin is smallest on the three suites held out for some version (\code{dlheld2}, \code{dlheld3}, \code{dlheld7})\\
\addlinespace
input dimension: $d=15,60,120$, 7 targets (\cref{sec:res-dscale}) & \code{nets}, run at every $d$ & exponent of $\nstar$ in $d$ 0.01--0.50 (parities of degree 2, 3, 4: 0.24, 0.30, 0.50), networks 0.91--2.40 (1.03, 1.96, 2.40); median factor $6.2\times$ at $d=15$, $42.5\times$ at $d=120$ & 7, 7, 4 development targets; 2 seeds at $d\ne30$; the weakest variant; 3 or 4 points per slope\\
\addlinespace
no search: \code{lowdeg}, 138 targets (\cref{sec:res-lowdeg}) & DFG alone; \code{nets_best} & \code{lowdeg} against DFG alone $+1.64$ over 122 ($3.1\times$ more data; worse on 112, better on 16, never on 12 that DFG learns); against \code{nets_best} $-0.16$ (53 better, 62 worse) & in-sample; Part 6's check, unchanged; faster only where DFG alone needs $\le133$ points\\
\addlinespace
tabular: 2 datasets, 9 cells (\cref{sec:res-tabular}) & \code{mlp_best} (hindsight); \code{hgb}; v3 & cells better / tie / worse: 0/8/1 against \code{mlp_best}, 0/9/0 against \code{hgb}, 0/8/1 against v3 & kr\_vs\_kp and mushroom, the only 2 of 26 datasets with $\pm1$-codable inputs and two classes; no real-valued data\\
\addlinespace
post hoc: the core, 138 targets (\cref{sec:core}) & DFG alone, same seeds; \code{nets_best} & against DFG alone $+0.07$ over 133 (28 lower, 26 higher; loses par5); against \code{nets_best} $-1.74$ (111/9), DFG alone $-1.88$ (117/5); \code{dlheld7} $-1.30$; median FLOPs 0.094 of DFG's & seeds 0, 1; defined after every result had been read: nothing is held out for it\\
\bottomrule
\end{tabularx}

\caption{The results at a glance \status{Numerical}. A ratio is $\log_2(\nstar_{\rm DFG}/\nstar_{\rm ref})$ (for \code{lowdeg} and the core, of theirs to the reference's), negative when the first needs less data; means are over the targets both learn. ``Clearly'': beyond $\pm0.25$, or only one of the two reaches 10\% error; on tabular data, errors that differ by more than 10\% and 2 s.e.\ (\cref{sec:metrics}). The s.d.\ and s.e.\ on \code{dlheld7} are over its 11 learned targets, as independent draws from the family. Only \code{dlheld7} is held out for DFG v1.3. Generated by \repo{paper_dfg/make_tables.py} (\texttt{tab\_headline}) from \repo{study/part11/RESULTS.json} and \repo{study/part11/interp/ablations/ablations.json} (the core's FLOPs).}
\label{tab:headline}
\end{table}

\paragraph{What was found.} \Cref{tab:headline} gives the comparisons.
\begin{itemize}
\item \textbf{Less data than the networks} \status{Numerical} (\cref{sec:res-dlheld7,sec:res-suites,sec:res-dscale}). On the held-out \code{dlheld7}, DFG needs clearly less data than the best network variant on 10 of 12 targets and more on none, about $2.4\times$ less on average (mean $\log_2$ ratio $-1.29$, s.e.\ 0.27 if the 11 learned targets are independent draws from the family). On five of the ten wins logistic regression or default gradient boosting is also clearly better than the networks; on the other five, all sum-structured, DFG needs about $4.3\times$ less data and neither no-search baseline comes close. On seven targets run at $d=15$--60 (four also at $d=120$), DFG's $\nstar$ grows as $d^{0.01}$ to $d^{0.50}$ and the networks' as $d^{0.91}$ to $d^{2.40}$: a trade of samples for FLOPs, since an exhaustive degree-$k$ search costs of order $n\binom dk$ FLOPs, so that within DFG's budget of one network run the global search can pass no degree-4 character from $d=179$ on \status{Proved, by the code} (\cref{sec:res-dscale,sec:disc-why}). v3 needs $1.14\times$ less data than DFG alone on average over the 134 targets both learn, mostly on small formulas and thresholds that it learns by $n=128$ (\cref{sec:res-v3}).
\item \textbf{What carries the results} \status{Numerical, then Heuristic} (\cref{sec:res-lowdeg,sec:ablations}). Not the characters and the test alone: the low-degree algorithm, which shares DFG's characters of degree $\le4$ and the form of its test and nothing else (no signed-sum blocks, residual rounds or refitted readout), is level with the networks. Removing only the sum blocks from DFG costs most ($+0.91$ in mean $\log_2\nstar$ over the 59 of 65 ablation targets both learn, and 3 targets lost), then the residual-scaled tests and the small-$n$ machinery. A core defined post hoc, which keeps the characters of degree $\le3$, the unweighted signed-sum blocks without the eigenvector seeds, the residual-scaled tests (with $K=3$, so 0.29 nats more lenient), least squares and the small-$n$ cross-validation, matches DFG on average ($+0.07$ in mean $\log_2\nstar$ over the 133 targets both learn) at a median 0.094 of its FLOPs, but moves 54 of 138 targets beyond the tie band, 28 lower and 26 higher, where the paired null moves 8 (\cref{sec:core}).
\item \textbf{What it learns} \status{Numerical, then Heuristic} (\cref{sec:interp-atlas}). Parities are their single character; symmetric targets one exact signed-sum block; halfspaces the linear block's degree-1 characters, then blocks that encode the threshold; a parity of two majorities first the folded blocks $|t_A+t_B|$ and $|t_A-t_B|$, since $\mathrm{sign}(|t_A+t_B|^2-|t_A-t_B|^2)=\mathrm{sign}(t_At_B)$; a sum with a parity or gate partner ends in weighted blocks with the partner as inner term; formulas and multiplexers are their characters, after transient blocks that encode a clause or a branch. Where DFG errs, features are missing, not weights: in 133 of the 148 analysed cells with test error above 0.1, no weights on the learned features fit the target.
\item \textbf{Two regimes} \status{Proved, Numerical, then Heuristic} (\cref{sec:interp-when,sec:interp-fail}). On parities the test sets the data cost: a proved criterion predicts the outcome in all 250 cells of the 2-, 3- and 4-parities with $n\ge32$, and moving the test's level moves $\nstar$ by the predicted amount. On signed-sum targets, where the grid can tell, the evidence points to the search: on the 6 targets that DFG learns a doubling or more after their true block would first pass its test, a stricter test, which raises the block's threshold by about a quarter of a doubling of $n$, moves $\nstar$ by at most 0.05. On the other 24 learned sum targets the grid cannot separate the test from the search. At the extreme, the residue classes mod 3 have $1.6\cdot10^{-6}$ to $2.1\cdot10^{-4}$ of their Fourier weight at degrees 1--2, where the block search starts, and their block never enters DFG's model, although it would pass the test from $n=128$.
\item \textbf{Against the networks' order} \status{Numerical, then Heuristic} (\cref{sec:interp-order}). Inside DFG's first round its test statistics are positively rank-correlated with the Part 5 networks' order on all 11 targets (median Kendall $\tau$ 0.54 at $n=16384$), but the orders invert where a large coefficient sits at high degree: the multiplexer's largest term is DFG's first character by round-0 $z$-score and the networks' last. On the 8 targets whose coefficients differ in size, DFG's $n$ order follows coefficient size more than degree on 5, and the networks' time order follows degree on 7: DFG tends to admit characters in order of detectability ($n_{\rm glob}\propto\ln(4K\binom d{|S|}/\delta)/\hat f(S)^2$: degree enters only through a logarithm, coefficient size through its inverse square), the networks to acquire them in order of construction (lower degree first, products after their factors).
\end{itemize}

\needspace{5\baselineskip}
\paragraph{Scope.}
\begin{itemize}
\item \emph{What is held out.} For DFG v1.3, only \code{dlheld7}: 12 targets drawn from the blind family of signed-sum targets and fixed in a commit before v1.3 was designed (\cref{sec:part11}). DFG v1.3 was developed on two other draws from the same family, so \code{dlheld7} tests it on fresh targets from a distribution its designer had seen.
\item \emph{Everything else is in-sample.} The other 10 suites were used in developing some version of DFG, or were held out only for an earlier version, and several of DFG's parts were added after failures on their targets (\cref{tab:status}). Their numbers, the ablations, the core and the analyses describe the frozen learner; they do not test it.
\item \emph{The object was chosen post hoc}, after Part 10's results had been read, v3's on \code{dlheld7} included; DFG alone's results there were new, bar a disclosed timing check on two targets (\cref{sec:part11}).
\item \emph{The references} are one network recipe, not tuned per target, and \code{nets_best} is a hindsight minimum over the 2 to 5 variants run on each suite (\cref{sec:refs}).
\item \emph{The domain is Boolean}: inputs uniform on $\cube$, mostly with $d=30$, and two tabular datasets written as $\pm1$; nothing here is about real-valued data or images (\cref{sec:res-tabular}).
\item \emph{Seeds and noise.} $\nstar$ is read from test errors averaged over 3 training sets per target and $n$ (2 on \code{random} and at $d\ne30$), with no confidence interval. DFG's $\nstar$ on two different single training sets differ by a mean $|\log_2$ ratio$|$ of 0.339 (23 of 65 targets beyond $\pm0.25$) and refits on the same training sets with another column order by 0.012, so the $\pm0.25$ band is smaller than the single-seed spread and per-target verdicts are not significance calls (\cref{sec:metrics}). The ablations use one or two seeds, the core two, and most analyses seed 0.
\item \emph{AI-written}: code, runs, analyses and text, checked by AI verifiers; no human has reviewed them (\cref{app:provenance}).
\end{itemize}
Claims are labelled \status{Proved} (a proof is given, here or in the theory notes, which an independent verifier checked), \status{Numerical} (measured here, with the source file named) or \status{Heuristic} (a reading or a back-of-envelope argument, not tested here).

\paragraph{Contributions.}
\begin{enumerate}
\item What DFG v1.3 is, six statements about it, all but one exact, and its relation to familiar learners (\cref{sec:learner}; every step in \cref{app:learner-full}, proofs and checks in \cref{app:proofs}).
\item A pre-registered evaluation against the networks, the no-search baselines, v3 and the low-degree algorithm, in the input dimension, on tabular data and in compute (\cref{sec:protocol,sec:results}; in full in \cref{app:eval,app:further}).
\item What its parts do: 12 pre-registered ablations against a post-hoc paired null, and a post-hoc core DFG with nine parts removed (\cref{sec:ablations}; \cref{app:abl-detail}).
\item A descriptive, independently verified account of what DFG learns: its models, when structure enters relative to its test, its failures, its false admissions and its order of admission against the networks' (\cref{sec:interp}; \cref{app:interp-detail}).
\item What kind of learner DFG turned out to be, why it needs less data here, where it fails and what to study next (\cref{sec:discussion}); how Part 11 was run and checked (\cref{app:provenance}).
\end{enumerate}
